\chapter{Nitrites (Urine)}

\section*{What Is This Test?}
A urinary nitrite test detects nitrites in the urine as a surrogate marker of bacteriuria. Many Gram-negative uropathogens (Escherichia coli, Klebsiella, Proteus, Enterobacter, Pseudomonas) reduce dietary nitrate to nitrite via bacterial nitrate reductase; nitrite therefore indicates significant bacterial colonization. The test is a routine dipstick pad that is most specific (highly predictive when positive) for urinary tract infection when combined with the leukocyte esterase pad. The trade-off is reduced sensitivity: a negative test does not exclude UTI. False negatives occur with empty bladder dwell time $<$ 4 hours, low dietary nitrate, and organisms that do not reduce nitrate (Gram-positive enterococci, Staphylococcus saprophyticus).

\section*{Alternative Names}
Urine Nitrite; Dipstick Nitrite; Nitrituria; Urinalysis Nitrite

\section*{Principle}
The dipstick pad uses Griess reaction: nitrite reacts with p-arsanilic acid at acidic pH to form a diazonium compound, which couples with N-(1-naphthyl)ethylenediamine to form a pink-red azo dye. Detection threshold: 0.06--0.15 mg/dL nitrite. Color intensity graded as negative or positive (no scaling as it is qualitative).

\section*{History}
The Griess reaction was first described in 1858 and adapted to urine dipsticks by the 1950s. Combined with the leukocyte esterase reaction, it became a cornerstone of UTI screening.

\section*{Why This Test Is Ordered}
\begin{itemize}
  \item Screening for urinary tract infection in symptomatic patients (dysuria, frequency, urgency)
  \item Routine urinalysis screening
  \item Evaluation of asymptomatic bacteriuria in pregnancy
  \item Pre-operative evaluation before urologic procedures
  \item Long-term care or catheter surveillance (interpret with caution, colonization common without infection)
\end{itemize}

\section*{Is This a Primary or Secondary Test}
Primary screening test; a positive result with symptoms warrants urine culture confirmation.

\section*{How This Test Is Performed}
A first-morning urine sample is preferred because it reflects $>$ 4 hours of bladder dwell time needed for bacteria to convert nitrate to nitrite. Midstream clean-catch is standard. Dipstick pad is read at 60 seconds. Positive result reflexes to urine culture and susceptibility testing.

\section*{What Preparation Is Needed}
None specific. An adequate bladder dwell time ($>$ 4 hours) is essential. Patients on low-nitrate diets (vegetarian, restricted diets) may have false negatives. Vitamin C can interfere.

\section*{Contraindications and Precautions}
False negatives: short bladder dwell time ($<$ 4 hours), low dietary nitrate, organisms that do not reduce nitrate (Enterococcus faecalis, Staphylococcus saprophyticus), high urinary vitamin C, frequent voiding, dilute urine. False positives: rare with contaminated specimens, pink-tinted medications (phenazopyridine), expired dipsticks. Phenazopyridine (Pyridium) can cause false positive.

\section*{Risks}
None.

\section*{What Happens After the Test}
A positive nitrite with symptoms $+$ leukocyte esterase positive strongly suggests UTI: empirical antibiotics with culture confirmation. Negative nitrite with symptoms does not exclude UTI; proceed with urinalysis microscopy and culture if clinical concern.

\section*{Understanding Results}

\subsection*{Below Normal}
Negative nitrite is normal. It does not exclude UTI: Gram-positive organisms, short bladder dwell time, or low nitrate intake all produce false negatives.

\subsection*{Above Normal}
\begin{itemize}
  \item Significant bacteriuria from nitrate-reducing Gram-negative organisms (\textit{E.\ coli} responsible for 80\% of community UTIs)
  \item Asymptomatic bacteriuria, especially in pregnancy or preoperative evaluation
  \item Catheter-associated asymptomatic bacteriuria in long-term catheterized patients
  \item False positives with medications (phenazopyridine) or stale contaminated specimens
\end{itemize}

\section*{Normal Results}
\begin{itemize}
  \item \textbf{Nitrite (urine dipstick):} Negative in normal urine
  \item \textbf{Interpretation when positive:} Strongly suggests bacteriuria if bladder dwell time adequate
  \item \textbf{Combined with leukocyte esterase:} Both positive increases specificity for UTI
  \item \textbf{Cutoff concerns:} Detection requires $>$ 10$^5$ CFU/mL in most laboratories
\end{itemize}

\section*{Follow-up Tests}
\begin{itemize}
  \item \textbf{Urine culture and susceptibility:} Required for definitive diagnosis and treatment
  \item \textbf{Urinalysis with microscopic examination:} Confirms leukocyte esterase, WBC, bacteria
  \item \textbf{Leukocyte esterase pad of dipstick:} Indirect detection of pyuria
  \item \textbf{Pregnancy test (if applicable):} If UTI suspected in women of childbearing age
  \item \textbf{Renal ultrasound:} For recurrent UTIs or suspected obstruction
\end{itemize}

\section*{References}
\begin{enumerate}
  \item Hooton TM, Gupta K. Acute uncomplicated cystitis and pyelonephritis in women. In: Bennett JE, et al., eds. Mandell, Douglas, and Bennett's Principles and Practice of Infectious Diseases. 9th ed. Elsevier; 2020.
  \item Devillé WLJM, et al. The urine dipstick test useful to rule out infections. A meta-analysis. BMC Fam Pract. 2004;5:4.
  \item Strasinger SK, Schumann GB. Urinalysis and Body Fluids. 7th ed. FA Davis; 2020.
  \item Nicolle LE, et al. Clinical practice guideline for the management of asymptomatic bacteriuria. Clin Infect Dis. 2019;68(10):1611-1616.
\end{enumerate}

\clearpage